\documentclass[11pt]{article}
\usepackage[margin=1in]{geometry}
\usepackage{amsmath,amssymb,amsthm}
\usepackage[hidelinks]{hyperref}
\newtheorem{theorem}{Theorem}
\newcommand{\R}{\mathbb R}
\newcommand{\Ric}{\operatorname{Ric}}
\newcommand{\Scal}{\operatorname{Scal}}
\newcommand{\tr}{\operatorname{tr}}
\newcommand{\Der}{\operatorname{Der}}
\newcommand{\norm}[1]{\lVert #1\rVert}
\title{Local Ricci-pinching maxima on an almost-abelian group}
\author{Alec Kriebel\\
\small ORCID: \href{https://orcid.org/0009-0001-9320-500X}{0009-0001-9320-500X}}
\date{Research note --- 13 October 2026}

\begin{document}
\maketitle
\begin{abstract}
Lauret's July 2018 talk announced nonglobal local maxima of the Ricci
pinching functional that are not solvsolitons of the specified transitive
group. This note gives a self-contained small-parameter proof for the
known Lauret--Will family \(A_u=J\oplus uE_{12}\).
For \(0<u<1/2\), its standard metric is a non-strict local maximum of
\(F=\Scal^2/\norm{\Ric}^2\), with value \(1/3\), among all nearby
left-invariant metrics on its fixed five-dimensional simply connected
group. The metric fails that group's equation \(\Ric=cI+D\),
\(D\in\Der(\mathfrak g)\), and a same-group metric has value \(9/17\).
The proof uses a full orthogonal slice and a uniform remainder estimate
along an exact constant-value family. The metric is nevertheless a Ricci
soliton, isometric to \(\R^2\times H_3\). No novelty, first assertion, or
first proof is claimed. This AI-assisted research note is unrefereed.
\end{abstract}

\section{Statement and source context}

The existence statement addressed here was already announced by Jorge
Lauret on 23 July 2018, in a displayed theorem attributed to
Lauret--Will \cite[slide 13/17, overlay PDF pp.~71--73]{talk}.
It asserts the existence of Ricci-soliton local maxima that are
nonglobal and are not solvsolitons. The checked slides do not supply a
construction or proof of that clause.
The exact family used below, denoted \(D_t\), and its value \(1/3\)
are in Lauret--Will \cite[\S6.4, p.~19]{LWv1} and
\cite[\S6.4, p.~20]{LWv2}. Their critical-point classification and
nonglobality results are also prior work
\cite[Theorem 1.5, Theorem 6.5(iii), Proposition 6.11]{LWv2}.

Both checked arXiv versions explicitly leave small nilpotent
coefficients unresolved immediately after Lemma 6.12
\cite[pp.~18--19]{LWv1}\cite[p.~19]{LWv2}.
These statements postdate the talk announcement. This note does not
infer the reason for their difference. The full journal text was not
retrieved, so the statements about passages and page numbers concern
the specified arXiv versions. The proof below stands on its own;
neither the source chronology nor a claim of present community-open
status is a premise.

Question 14 in \S11 of the Morgan--Pansu collection
\cite[p.~12]{MP}, proposed by Lauret, asks whether local maxima of \(F\)
among left-invariant metrics on a fixed solvable group are
solvsolitons. In that question a solvsoliton means
\begin{equation}\label{eq:solvsoliton}
  \Ric_g=cI+D,\qquad c\in\R,\quad D\in\Der(\mathfrak g),
\end{equation}
where \(\mathfrak g\) is the Lie algebra of the specified group.
Thus the group presentation is part of the assertion.
Here and throughout, \(\Ric_g\) denotes the Ricci endomorphism and its
norm is the Hilbert--Schmidt norm with respect to \(g\).
Local maximality uses the usual topology on the cone of positive
definite inner products on \(\mathfrak g\).

Let \(J=\begin{pmatrix}0&-1\\1&0\end{pmatrix}\), and let \(E_{12}\)
be the \(2\) by \(2\) matrix with its only nonzero entry in position
\((1,2)\). For \(u>0\), define
\[
 A_u=\begin{pmatrix}
 0&-1&0&0\\1&0&0&0\\0&0&0&u\\0&0&0&0
 \end{pmatrix}=J\oplus uE_{12}.
\]
Let \(\mathfrak g_u=\R e_0\oplus V\), with
\(V=\operatorname{span}(e_1,e_2,e_3,e_4)\), have bracket
\([e_0,X]=A_uX\) and \([V,V]=0\).
This is a solvable Lie algebra: its derived algebra is contained in
the abelian ideal \(V\). Write \(G_u\) for its simply connected group
and \(g_u\) for the metric making the displayed basis orthonormal.

\begin{theorem}\label{thm:main}
For each fixed \(0<u<1/2\), the metric \(g_u\) is a non-strict local
maximum of \(F\) among all left-invariant metrics on \(G_u\), and
\(F(g_u)=1/3\). It does not satisfy \eqref{eq:solvsoliton} on
\(\mathfrak g_u\). It is not a global maximum: a left-invariant metric
on the same group has \(F=9/17\).
\end{theorem}

This supplies a negative answer to the literal fixed-group assertion
in Question 14, already announced in the talk cited above.
It is compatible with the Ricci-soliton local-maximum conclusion of
Lauret--Will \cite[Theorem 1.5]{LWv2}.
Section \ref{sec:presentation} explains the metric's Ricci-soliton
presentation explicitly. No conclusion for completely solvable groups
or groups of real type follows from this example.

\section{Curvature and the space of metrics}

We first make the curvature computation and the all-metrics reduction
explicit. For an arbitrary real matrix \(A:V\to V\), consider the
orthonormal almost-abelian presentation \([e_0,X]=AX\). Put
\[
 S=\tfrac12(A+A^T),\qquad K=\tfrac12(A-A^T).
\]
For left-invariant \(X,Y\in V\), the Koszul formula gives
\begin{equation}\label{eq:connection}
\begin{aligned}
 \nabla_{e_0}e_0&=0,& \nabla_{e_0}X&=KX,\\
 \nabla_Xe_0&=-SX,& \nabla_XY&=\langle SX,Y\rangle e_0.
\end{aligned}
\end{equation}
Use the convention
\(R(U,W)=\nabla_U\nabla_W-\nabla_W\nabla_U-\nabla_{[U,W]}\).
Substituting \eqref{eq:connection} yields
\[
 R(e_0,X)e_0=(S^2+[S,K])X,\qquad
 R(X,Y)Z=-\langle SY,Z\rangle SX+\langle SX,Z\rangle SY.
\]
Taking the Ricci trace gives
\begin{equation}\label{eq:ricci}
 \Ric_{00}=-\tr S^2,\qquad
 \Ric_{0V}=0,\qquad
 \Ric|_V=[K,S]-(\tr A)S
       =\tfrac12[A,A^T]-(\tr A)S.
\end{equation}
For example, the \(e_0\) contribution to \(\Ric|_V\) is
\(-S^2-[S,K]\), and the trace over \(V\) is
\(S^2-(\tr S)S\), giving the stated expression.
This also recovers the almost-abelian formula in
\cite[equation (3)]{LWv2}.

Every matrix used below has trace zero. With the Frobenius norm, put
\[
 s(A)=\norm{S}^2,\qquad q(A)=\norm{[A,A^T]}^2.
\]
Equation \eqref{eq:ricci} implies
\begin{equation}\label{eq:reduction}
 \Scal=-s,\qquad \norm{\Ric}^2=s^2+\tfrac14q,\qquad
 F(A)=\frac{s^2}{s^2+q/4}.
\end{equation}
When \(s>0\), this reduces the desired inequality to
\begin{equation}\label{eq:P}
 F(A)\le\tfrac13
 \quad\Longleftrightarrow\quad
 P(A):=q(A)-8s(A)^2\ge0.
\end{equation}
For \(A_u\), direct multiplication gives
\(s(A_u)=u^2/2\), \(q(A_u)=2u^4\), and hence \(F(A_u)=1/3\).
The function \(F\) is invariant under nonzero scalar multiplication
of \(A\) and under orthogonal conjugation.

Fix \(u=u_0\) in the theorem's range. A nearby inner product on the
fixed vector space \(\R e_0\oplus V\) has positive definite Gram matrix
\[
 \Gamma=\begin{pmatrix}\alpha&\beta^T\\\beta&\Gamma_V\end{pmatrix}.
\]
Set
\[
 T=\Gamma_V^{-1/2},\qquad
 c=(\alpha-\beta^T\Gamma_V^{-1}\beta)^{-1/2}.
\]
The columns of \(T\) give an orthonormal basis of \(V\), and
\(f_0=c(e_0-\Gamma_V^{-1}\beta)\) is its unit orthogonal complement.
The component of \(f_0\) in \(V\) has zero adjoint action on \(V\).
Thus the bracket matrix in this orthonormal basis is exactly
\begin{equation}\label{eq:allmetrics}
 cT^{-1}A_{u_0}T.
\end{equation}
As \(\Gamma\to I\), both \(c\to1\) and \(T\to I\) smoothly.
Consequently it suffices to prove \eqref{eq:P} for every conjugate of
\(A_{u_0}\) sufficiently close to \(A_{u_0}\).
Conversely, any conjugate is realized by a metric on the fixed group:
an arbitrary invertible change of basis in \(V\), declared orthonormal
with \(e_0\), gives such a matrix. Cross terms in \(\Gamma\) have
therefore been fully included in the local reduction.

\section{A full orthogonal slice}

Let \(A\) be a nearby conjugate of \(A_{u_0}\).
Its generalized zero eigenspace
\[
 W=\ker A^2
\]
is two-dimensional and invariant. The relation \(A^4+A^2=0\) implies
that \(E=I+A^2\) is an idempotent with image \(W\).
This is a smooth projection, generally not an orthogonal one.
We can construct a smooth local orthonormal frame directly:
put \(p=Ee_4\) and
\[
 f_3=\frac{Ap}{\norm{Ap}},\qquad
 f_4=\frac{p-\langle p,f_3\rangle f_3}
 {\norm{p-\langle p,f_3\rangle f_3}}.
\]
At \(A_{u_0}\) the denominators are \(u_0\) and \(1\), respectively,
so neither vanishes nearby. Since \(A^2p=0\), these choices give
\[
 Af_3=0,\qquad Af_4=vf_3,\qquad
 v=\frac{\norm{Ap}}{\norm{p-\langle p,f_3\rangle f_3}}>0,
 \quad v\longrightarrow u_0.
\]
Project \(e_1,e_2\) orthogonally onto \(W^\perp\) and apply
Gram--Schmidt to obtain a smooth frame \(f_1,f_2\) close to \(e_1,e_2\).

In the frame \(f_1,f_2,f_3,f_4\), invariance of \(W\) gives
\begin{equation}\label{eq:slice}
 M=\begin{pmatrix}C&0\\L&N\end{pmatrix}
  =\begin{pmatrix}
  a&b-k&0&0\\ b+k&-a&0&0\\
  x&y&0&v\\ z&w&0&0
  \end{pmatrix},
 \qquad k=\sqrt{1+a^2+b^2}.
\end{equation}
Indeed, \(C\) represents the quotient action on \(V/W\), which has
eigenvalues \(i,-i\). It has trace zero and determinant one.
Writing its symmetric part as
\(\begin{pmatrix}a&b\\b&-a\end{pmatrix}\) and skew part as \(kJ\)
gives \(k^2-a^2-b^2=1\); proximity to \(J\) chooses \(k>0\).
All six transverse coordinates
\(\xi=(a,b,x,y,z,w)\) tend to zero as \(A\to A_{u_0}\).
The lower-left block \(L\) is unrestricted in the estimates that follow.
This point is essential: \(W^\perp\) need not be invariant.

For a checkable expansion, block multiplication gives
\[
 [M,M^T]=
 \begin{pmatrix}
 [C,C^T]-L^TL&CL^T-L^TN\\
 LC^T-N^TL&LL^T+[N,N^T]
 \end{pmatrix}.
\]
It follows that
\begin{align}
 q(M)&=\norm{[C,C^T]-L^TL}^2+
 2\norm{LC^T-N^TL}^2+\norm{LL^T+[N,N^T]}^2,\label{eq:blockq}\\
 s(M)&=2(a^2+b^2)+\tfrac12(x^2+y^2+z^2+w^2+v^2).\label{eq:blocks}
\end{align}
In particular, \(s(M)\ge v^2/2>0\).
Put \(h^2=a^2+b^2\), \(r^2=x^2+y^2+z^2+w^2\), and \(d=wx-yz\).
Expanding \eqref{eq:blockq}--\eqref{eq:blocks}, using \(k^2=1+h^2\),
gives the exact identity
\begin{equation}\label{eq:QR}
 P(M)=Q_v(\xi)+R_v(\xi),
\end{equation}
where
\begin{equation}\label{eq:Q}
\begin{aligned}
 Q_v(\xi)={}&16(2-v^2)(a^2+b^2)
       +2(x+vw)^2+2(y-vz)^2\\
       &+2(1-4v^2)(z^2+w^2)
\end{aligned}
\end{equation}
and
\begin{equation}\label{eq:R}
\begin{aligned}
 R_v(\xi)={}&-12h^2r^2-24ak(wz+xy)
       +12bk(x^2+z^2-y^2-w^2)\\
 &+4v\{a(wy-xz)-b(wx+yz)\}
       +4v(k-1)d-4d^2.
\end{aligned}
\end{equation}
These formulas exhibit \(P(v,0)=0\) and \(D_\xi P(v,0)=0\) for
every \(v>0\). Each term in \(R_v\) has order at least three in
\(\xi\), because \(k-1=h^2/(k+1)\).

\section{Uniform positivity and local maximality}

Choose a closed interval \(I=[v_-,v_+]\) with
\(0<v_-<u_0<v_+<1/2\), and set \(c_0=1-4v_+^2>0\).
The two mixed-variable quadratic blocks in \eqref{eq:Q}, before the
factor \(2\), have matrices
\[
 \begin{pmatrix}1&v\\v&1-3v^2\end{pmatrix},
 \qquad
 \begin{pmatrix}1&-v\\-v&1-3v^2\end{pmatrix}.
\]
Their determinants are \(1-4v^2\) and their traces are \(2-3v^2\).
For a positive definite \(2\) by \(2\) matrix, its smaller eigenvalue
is at least its determinant divided by its trace. Thus the smaller
eigenvalue after the factor \(2\) is at least
\[
 \frac{2(1-4v^2)}{2-3v^2}\ge 1-4v^2\ge c_0.
\]
The coefficient \(16(2-v^2)\) also exceeds \(c_0\). Hence, with
\(\rho=\norm{\xi}\),
\begin{equation}\label{eq:Qbound}
 Q_v(\xi)\ge c_0\rho^2\qquad(v\in I).
\end{equation}

For completeness, an explicit uniform bound on \eqref{eq:R} suffices.
If \(\rho\le1\), then \(k\le\sqrt2\), \(k-1\le h^2/2\), and
\[
 |wz+xy|,\ |wy-xz|,\ |wx+yz|,\ |d|\le r^2/2,\qquad
 |x^2+z^2-y^2-w^2|\le r^2.
\]
Using \(v<1/2\) and \(|a|+|b|\le\sqrt2\rho\), the six terms of
\eqref{eq:R}, in their displayed order, are bounded respectively by
\[
 12\rho^4,\quad
 12\sqrt2\rho^3,\quad
 12\sqrt2\rho^3,\quad
 \sqrt2\rho^3,\quad
 \tfrac12\rho^4,\quad
 \rho^4.
\]
Consequently
\begin{equation}\label{eq:Rbound}
 |R_v(\xi)|\le(\tfrac{27}{2}+25\sqrt2)\rho^3<50\rho^3
 \qquad(v\in I,\ \rho\le1).
\end{equation}
For \(\rho\le\min(1,c_0/100)\), equations
\eqref{eq:QR}, \eqref{eq:Qbound}, and \eqref{eq:Rbound} give
\[
 P(M)\ge c_0\rho^2-50\rho^3\ge\tfrac12c_0\rho^2\ge0.
\]
Continuity of the slice coordinates ensures that every sufficiently
nearby conjugate has \(v\in I\) and satisfies this bound on \(\rho\).
Orthogonal invariance, \eqref{eq:allmetrics}, and \eqref{eq:P}
therefore give \(F(g)\le1/3\) for every sufficiently nearby
left-invariant metric on \(G_{u_0}\).

The estimate is uniform along the flat parameter \(v\); this is why
the argument proves a neighborhood maximum despite a degenerate
Hessian. The family \(A_v\) consists of conjugates of \(A_{u_0}\)
by \(\operatorname{diag}(1,1,v/u_0,1)\) for all \(v>0\),
and \(F(A_v)=1/3\). In particular the maximum is
non-strict. The theorem asserts only the open range \(0<u<1/2\);
neither endpoint is included.

\section{The fixed-group obstruction and a larger value}

At the candidate, \eqref{eq:ricci} gives, in the order
\(e_0,e_1,e_2,e_3,e_4\),
\begin{equation}\label{eq:riccandidate}
 \Ric=\operatorname{diag}(-u_0^2/2,0,0,u_0^2/2,-u_0^2/2).
\end{equation}
If \(\Ric=cI+D\), the required map \(D=\Ric-cI\) is this specific
diagonal map. Write its entries \(d_0,\ldots,d_4\).
The derivation identity for \([e_0,e_1]=e_2\) requires
\(d_2=d_0+d_1\). Since \(d_1=d_2=-c\), it forces \(d_0=0\)
and \(c=-u_0^2/2\). Then \(d_3=u_0^2\) and \(d_4=0\).
The nonzero bracket \([e_0,e_4]=u_0e_3\) instead requires
\(d_3=d_0+d_4=0\), a contradiction.
This proves failure of the precise equation on \(\mathfrak g_{u_0}\).

To check nonglobality directly, set
\[
 B=\begin{pmatrix}
 0&-1&0&0\\1&0&0&0\\1&0&0&1\\0&-1&0&0
 \end{pmatrix},\qquad
 T_B=\begin{pmatrix}I_2&0\\Y&I_2\end{pmatrix},
 \qquad Y=\begin{pmatrix}0&2\\1&0\end{pmatrix}.
\]
With \(H=\operatorname{diag}(1,1,1/u_0,1)\), multiplication gives
\[
 HA_{u_0}H^{-1}=A_1,\qquad
 T_BA_1T_B^{-1}=B.
\]
Thus \(B\) represents a metric on the same fixed group.
Its invariants are \(s(B)=3/2\) and \(q(B)=8\), so
\[
 F(B)=\frac{(3/2)^2}{(3/2)^2+8/4}=\frac9{17}>\frac13.
\]
This proves the last assertion of Theorem \ref{thm:main}.
Nonglobality for the known family also follows from the prior
optimization result \cite[Theorem 6.5(iii)]{LWv2};
the displayed comparison is an explicit check independent of that theorem.

\section{The Ricci-soliton presentation and disclosure}
\label{sec:presentation}

The obstruction above depends on the transitive group.
This distinction can be seen without a classification theorem.
Identify the simply connected \(G_{u_0}\) with \(\R\times\R^4\),
with multiplication
\[
 (t,X)(s,Y)=(t+s,X+e^{tA_{u_0}}Y).
\]
Writing \(X=(x_1,x_2,x_3,x_4)\), its standard left-invariant metric is
\[
 dt^2+\norm{e^{-tA_{u_0}}dX}^2
 =dx_1^2+dx_2^2+dt^2+(dx_3-u_0t\,dx_4)^2+dx_4^2.
\]
The rotation block contributes the Euclidean factor
\(dx_1^2+dx_2^2\). The remaining three terms are the
standard left-invariant metric on the Heisenberg group \(H_3\)
whose orthonormal basis \(T,Z,W\) satisfies
\([T,W]=u_0Z\). Hence the Riemannian manifold is the product
\(\R^2\times H_3\).

On this product Lie algebra, in the order of the Euclidean basis,
\(T,Z,W\), the Ricci operator is
\(\operatorname{diag}(0,0,-u_0^2/2,u_0^2/2,-u_0^2/2)\).
It equals \(cI+D_{\mathrm{prod}}\), where
\[
 c=-3u_0^2/2,\qquad
 D_{\mathrm{prod}}=
 \operatorname{diag}(3u_0^2/2,3u_0^2/2,u_0^2,2u_0^2,u_0^2).
\]
This map is a derivation: the central entry is the sum of the two
noncentral Heisenberg entries, and the Euclidean summand is abelian.
The product presentation is therefore a solvsoliton and, consequently,
a Ricci soliton: the symmetric derivation integrates to automorphisms,
whose pullback metric has derivative \(2g(D_{\mathrm{prod}}\,\cdot,\cdot)\).
Thus its Ricci tensor is a scalar multiple of the metric plus a Lie
derivative. Theorem \ref{thm:main} does not assert the existence of
a local maximum that is not a Ricci soliton.

\paragraph{Status, credit, and AI assistance.}
This is an unrefereed research note. The known family,
value, critical-point restrictions, and nonglobality machinery are
credited above; the earlier exact existence announcement is explicit
in the abstract and introduction. No new-family claim, firstness
claim, established novelty, or current community-open claim is made.
The direct proof presented here is intended to be checked from the
displayed identities and estimates.
AI tools were used extensively in solving the problem, symbolic
computation, source review, writing, and adversarial verification.
Machine agents are not independent human referees; this work has not
received external human peer review. The author retains responsibility
for the claims, calculations, citations, and exposition.

\begin{thebibliography}{9}
\bibitem{MP}
F.~Morgan and P.~Pansu,
\emph{A list of open problems in Differential Geometry},
collection associated with Modern Trends in Differential Geometry,
S\~ao Paulo, July 2018, \S11, Question 14, p.~12.
\href{https://www.imo.universite-paris-saclay.fr/~pierre.pansu/problems_MTDG.pdf}
{Accessible primary PDF}. Accessed 13 October 2026.

\bibitem{talk}
J.~Lauret,
\emph{Homogeneous Ricci curvature and the beta operator},
S\~ao Paulo talk, 23 July 2018, slide 13/17,
overlay PDF pp.~71--73.
\href{https://www.ime.usp.br/~mtg/slides/j_lauret.pdf}{Primary slides}.
Accessed 13 October 2026.

\bibitem{LWv1}
J.~Lauret and C.~E.~Will,
\emph{The Ricci pinching functional on solvmanifolds},
\href{https://arxiv.org/abs/1808.01380v1}{arXiv:1808.01380v1},
3 August 2018. Version-specific references in this note use its
printed pagination.

\bibitem{LWv2}
J.~Lauret and C.~E.~Will,
\emph{The Ricci pinching functional on solvmanifolds},
\href{https://arxiv.org/abs/1808.01380v2}{arXiv:1808.01380v2},
9 May 2019. Journal metadata: \emph{Quarterly Journal of Mathematics}
\textbf{70} (2019), 1281--1304,
\href{https://doi.org/10.1093/qmath/haz020}{doi:10.1093/qmath/haz020}.
The full text checked here is arXiv v2; no assertion of
textual identity with the journal version is made.
\end{thebibliography}
\end{document}
